\chapter{Health}\label{chap:health}

This chapter collects the parts of the model that concern health: a social
health insurance fund, the efficiency of factors that rises with health
services, the days of work lost to illness, and the death rate, which changes
population and labour in dynamic runs.

\section{A social health insurance fund}\label{sec:fund}

A fund mapped with \manref{cge_map}, option \cmd{insurance()}, is financed by
contributions at the rate $\V{tshi}_f$ on the income of each factor, taken
before the income is distributed \eqref{eq:I2}, and by transfers from the
government (in real terms) and the rest of the world (in foreign currency)
\eqref{eq:I12}. It buys a fixed basket of health services,
$\V{QSHI}_c = \overline{\V{qshi}}_c\,\V{SHIADJ}$ \eqref{eq:I13}, scaled so
that it spends its revenue less a fixed share saved \eqref{eq:I14}. A higher
contribution rate or larger transfers therefore raise its purchases. Its
savings enter the savings--investment balance \eqref{eq:S5}. The calibration
of the fund is given in section~\ref{sec:cal-institutions}; its parameters
can be shocked (\cmd{tshi}, \cmd{qshi}, \cmd{trshi}, \cmd{trshir}). A
donor-funded non-profit institution is mapped as a household that receives
transfers from the rest of the world.

\section{Health services and factor efficiency}\label{sec:efficiency}

Health outcomes respond to health services per head relative to the base,
\begin{equation*}
R = \V{HS} / \V{HS}^0,
\end{equation*}
where $\V{HS}$ is the volume of the health services in $HC$ bought by
households, the government and the health insurance fund \eqref{eq:H1}, and
$\V{HS}^0$ its base value, which grows with the population in dynamic runs
(chapter~\ref{chap:dynamics}). Each outcome has a constant elasticity with
respect to $R$. These elasticities are hypotheses about how services become
health, to be taken from studies of the country or region (for example
\citealt{bokhari2007}); the package supplies none. An outcome whose elasticity
is zero is constant unless shocked, and needs no health commodities.

The efficiency of factor $f$ is $\V{EFF}_f = R^{\eta_f}$ \eqref{eq:H2}, with
the elasticity $\eta_f$ given by \manref{cge_calibrate}, options
\cmd{healthcom()}, \cmd{healthfac()} and \cmd{healthelas()}. It multiplies the
factor's input in value added (section~\ref{sec:production}). The exogenous
efficiency $\V{af}_f$ can be shocked directly.

\section{Days lost to illness}\label{sec:sickdays}

For a factor $f \in FS$ the user gives the base share of working time lost to
illness, $0 \le s^0_f < 1$ (\cmd{sickdays()}), for example from a labour force
or household survey, and its elasticity $\epsilon^s_f$ with respect to $R$
(\cmd{sickelas()}; negative when more services mean fewer days lost). The
share lost is $S_f = \sigma_f R^{\epsilon^s_f}$, with $\sigma_f = s^0_f$ unless
shocked (\cmd{sick}, for example an epidemic), and the days worked per worker,
relative to the base, are
\begin{equation}
\V{DW}_f = \frac{1 - \sigma_f R^{\epsilon^s_f}}{1 - s^0_f}, \qquad f \in FS
\tag{H3}\label{eq:H3}
\end{equation}
and $\V{DW}_f = 1$ otherwise. Days worked scale the factor's effective input
in production, $\widetilde{\V{QF}}_{fa} = \V{af}_f\,\V{EFF}_f\,\V{DW}_f\,\V{QF}_{fa}$
(\eqref{eq:Q3}, \eqref{eq:Q4}). Factor quantities $\V{QF}_{fa}$ and
$\V{QFS}_f$, the price $\V{W}_f$ and unemployment (chapter~\ref{chap:labour})
stay per worker. In real terms fewer days per worker are the same as less
labour measured in days: with $L$ workers and $D$ days each, a wage $w$ per
effective unit pays $w\,D\,L$ whether the factor is counted in workers or in
days. $\V{EFF}_f$ measures output per day worked, and $\V{DW}_f$ the days
worked. $S_f$ must stay below one.

\section{Mortality}\label{sec:mortality}

The user may give the base crude death rate $m^0$, deaths per head per period
(\cmd{mortality()}), and its elasticity $\epsilon^m$ with respect to $R$
(\cmd{mortelas()}). The death rate is
\begin{equation*}
\V{MR} = \mu^m R^{\epsilon^m},
\end{equation*}
with $\mu^m = m^0$ unless shocked (\cmd{mort}). It is reported, and does not
enter the equations of a period: deaths in period $t$ change the population
and labour from period $t + 1$ (section~\ref{sec:population}), so each period
is solved with its population given.

\section{Population and labour in dynamic runs}\label{sec:population}

The population growth rate $g_P$ given for a run (\manref{cge_dynamic},
option \cmd{population()}) is taken to hold at the base death rate. With a
death rate, the population grows from $t - 1$ to $t$ at
\begin{equation*}
g_{P,t} = g_P + \bigl(m^0 - \V{MR}_{t-1}\bigr),
\end{equation*}
which scales the subsistence quantities, $\V{HS}^0$ and the population index
$\V{POP}$ (chapter~\ref{chap:dynamics}). For the factors named in
\cmd{deathfactors()} (labour, not capital), deaths at working ages change the
supply:
\begin{equation*}
g_{f,t} = g_f + \frac{\theta^D}{\alpha^W}\bigl(m^0 - \V{MR}_{t-1}\bigr),
\end{equation*}
where $\theta^D$ is the share of deaths at working ages (\cmd{deathshare()})
and $\alpha^W$ the working-age share of the population (\cmd{workingage()}):
excess deaths per head of the population, $\V{MR} - m^0$, remove a share
$\theta^D(\V{MR} - m^0)/\alpha^W$ of the working-age population, and the
factor's workers in the same proportion. The labour force of such a factor
(chapter~\ref{chap:labour}) grows at the same rate. $\V{MR}_{t-1}$ uses the
parameters and health services of period $t - 1$, including its shocks.

In a baseline, rising health services per head lower the death rate and raise
population growth above $g_P$; scenarios are therefore compared with a
baseline run with the same settings.

\section{Reporting}

The health table of \manref{cge_report} shows, when the model has them:
health services per head $R$; factor efficiency; the share of working time
lost, $100\,S_f$ per cent, and the days worked $\V{DW}_f$ for $f \in FS$; the
death rate per 1,000, $1000\,\V{MR}$; the population index $\V{POP}$; and the
revenue and purchases of the insurance fund. Shares and rates change in
points.
